\documentclass[11pt]{article}
\usepackage[a4paper,margin=24mm]{geometry}
\usepackage{amsmath,amssymb,amsthm,booktabs,hyperref,fancyvrb}
\hypersetup{colorlinks=true,urlcolor=blue,linkcolor=blue}
\newtheorem{theorem}{Theorem}
\newtheorem{lemma}{Lemma}
\title{Three counterexamples to the
Akbari--Cameron--Khosrovshahi row-space conjecture}
\author{
Alex Chan\thanks{Harvard University and NBER.
Correspondence:
\href{mailto:achan@hbs.edu}{\texttt{achan@hbs.edu}}.}
\and
Weilin Meng\thanks{Harvard Business School.}
\and
Kurt Sweat\thanks{UT Southwestern Medical Center.}
}
\date{22 October 2026}

\hypersetup{pdfauthor={Alex Chan, Weilin Meng, Kurt Sweat}}
\begin{document}
\maketitle
\begin{abstract}
We give three connected simple graphs on 52 vertices whose adjacency
matrices have real rank 14 and whose only Boolean row-space vectors
are zero and the adjacency rows. Their different edge counts make
them mutually nonisomorphic. Each graph refutes the universal
Akbari--Cameron--Khosrovshahi adjacency-row-space conjecture.
For all three graphs, Lean 4 verifies the simple-adjacency conditions,
the existence of an edge, the classification of all Boolean vectors
in the real row space, and the resulting negation of the conjecture.
The finite classification checks all $2^{14}=16,384$ Boolean prefixes
using exact integer certificates. We also give a mathematical
blowup argument producing connected counterexamples of every order
at least 52.
\end{abstract}
\section*{Statement and result}
For a real matrix $A$, let $\operatorname{Row}_{\mathbb R}(A)$
denote its real row space. A Boolean vector has every coordinate
in $\{0,1\}$.

The universal ACK adjacency-row-space conjecture asserts that every
finite simple graph with at least one edge has a nonzero Boolean
vector in the real row space of its adjacency matrix that is not
an adjacency row. The originating Question~2 of
Akbari, Cameron, and Khosrovshahi~\cite{ack} asks which graphs
have this property. The universal assertion is explicitly stated
as Conjecture~1.1 by Akansha and Sivakumar~\cite{akansha}.
Sciriha et al.~\cite{sciriha} study restrictions on potential
counterexamples.
\begin{theorem}
For each of the explicitly specified graphs $G_{550},G_{563},G_{570}$ below,
\[
\operatorname{Row}_{\mathbb R}(A_G)\cap\{0,1\}^{52}
=\{\mathbf 0\}\cup\{\text{rows of }A_G\}.
\]
Each graph is connected, simple, and reduced, and has rank 14. In particular, the ACK universal row-space conjecture is false.
\end{theorem}
\begin{center}
\begin{tabular}{lrrrr}
\toprule
Graph & Vertices & Edges & Real rank & Boolean vectors\\
\midrule
$G_{550}$ & 52 & 550 & 14 & 53\\
$G_{563}$ & 52 & 563 & 14 & 53\\
$G_{570}$ & 52 & 570 & 14 & 53\\
\bottomrule
\end{tabular}
\end{center}
The Boolean-vector count includes zero. All 52 adjacency rows are distinct and nonzero. The different edge counts prove the graphs are mutually nonisomorphic.
\section*{Scope}
The three examples refute the same universal adjacency-row-space
assertion. They do not settle the separate maximum-order or
extremal-classification conjectures for reduced graphs of a given
rank. We make no minimality claim for order 52 or rank 14.

The formal proof sources are available at
\begin{center}
\url{https://github.com/alexgocardinal/ack-rowspace-counterexample}.
\end{center}
The verified proof revision is
\href{https://github.com/alexgocardinal/ack-rowspace-counterexample/tree/5b04adb}
{\texttt{5b04adb}}.
The repository is currently private; access is required.
\newpage
\section*{Exact construction}
All mask coordinates below are indexed $0,\ldots,13$: coordinate $j$ is bit $j$, so the least-significant bit is coordinate 0. Let $b$ be the following list of 14 masks:
\begin{Verbatim}[fontsize=\small]
4018, 11465, 5136, 6034, 8429, 721, 3634, 7995,
14985, 1513, 6863, 13763, 3468, 2322
\end{Verbatim}
Define the $14\times14$ matrix $B$ by $B_{ij}=\operatorname{bit}_j(b_i)$.
It is symmetric with zero diagonal, and $\det(B)=-9$.
For each graph take its following list $c$ of 38 masks, and let $W$ be the $52\times14$ Boolean matrix whose row masks are the concatenation $(b,c)$. Define
\[
                         A=W B^{-1} W^T.
\]
This formula, involving only the integers displayed here, uniquely specifies each graph. Exact arithmetic verifies that every entry of $A$ is 0 or 1 and every diagonal entry is zero. The first 14 rows and columns of $A$ form $B$.
\paragraph{$G_{550}$.}
\begin{Verbatim}[fontsize=\small]
384, 676, 1497, 1517, 1810, 2031, 2444, 2692,
2859, 2998, 3080, 3362, 3464, 4289, 4559, 4803,
5378, 5520, 5650, 5715, 6099, 6162, 6227, 6284,
6349, 6566, 8385, 9232, 9248, 9904, 10761, 12091,
13379, 13807, 14035, 14051, 14635, 15410
\end{Verbatim}
\paragraph{$G_{563}$.}
\begin{Verbatim}[fontsize=\small]
384, 449, 793, 1129, 1296, 1452, 1810, 2031,
2060, 2692, 3080, 3390, 3451, 4875, 5378, 5520,
6099, 6566, 6611, 6665, 6838, 7586, 8385, 8745,
9232, 9616, 12739, 13698, 13807, 13881, 14107, 14123,
14563, 15488, 963, 5650, 6349, 10445
\end{Verbatim}
\paragraph{$G_{570}$.}
\begin{Verbatim}[fontsize=\small]
89, 384, 493, 963, 1497, 2617, 2843, 3120,
3346, 3362, 3712, 4112, 5152, 5241, 5467, 6099,
6182, 6271, 6284, 8385, 8457, 9011, 9089, 9248,
9632, 10278, 10445, 10662, 12169, 12355, 13011, 13537,
13609, 13807, 14051, 14883, 14989, 15488
\end{Verbatim}
The repository supplies the full adjacency matrices as lists of integer row masks in the Lean witness files. In particular, the formulas specify ordinary simple unweighted graphs; the intermediate rational inverse does not define weighted edges.
\newpage

\section*{Proof and exact certificates}
Write $W=\binom{B}{U}$, with $U$ a $38\times14$ Boolean matrix. Then
\[
 A=\begin{pmatrix}B&U^T\\U&UB^{-1}U^T\end{pmatrix}.
\]
Its Schur complement is zero, so $\operatorname{rank}_{\mathbb R}(A)=14$.
More explicitly, every row-space vector has the unique form
\[
                 f(x)=(x,\ xB^{-1}U^T),\qquad x\in\mathbb R^{14}.
\]
Indeed, the first 14 rows of $A$ span its row space, and their nonsingular first block is $B$. A Boolean row-space vector must therefore have $x\in\{0,1\}^{14}$. There are exactly $2^{14}=16,384$ assignments to check.
\par\medskip
For each graph, exact rational arithmetic evaluates $xB^{-1}U^T$ for all these assignments. Exactly 53 assignments produce a Boolean vector. Their prefixes are zero and the 52 row masks of $W$; their full vectors are zero and the adjacency rows of $A$. The remaining 16,331 assignments each have a coordinate outside $\{0,1\}$. This proves the theorem.
\par\medskip
The displayed construction also permits direct checks of the
edge counts, connectivity, and distinctness of the adjacency rows.
The nonsingularity of $B$ and the block identity above prove the
rank assertion without a floating-point rank computation.

\section*{Lean formalization}
The repository formalizes the conjecture over $\mathbb R$.
Its predicate \texttt{InRealRowSpace} explicitly requires a vector
to be a real linear combination of adjacency rows.
The predicate \texttt{IsSimpleAdjacency} requires zero diagonal,
symmetry, and entries in $\{0,1\}$.
The conjecture also explicitly requires an edge and excludes both
the zero vector and every existing adjacency row.

For each graph, an integer matrix $D$ is checked to satisfy
\[
  3A_{ij}=\sum_{k=0}^{13} A_{ik}D_{kj}.
\]
Linearity then gives the same reconstruction for every vector
in the real row space. If its first 14 coordinates have mask $m$,
its reconstructed coordinate $j$ is
\[
  \frac{1}{3}\sum_{\substack{0\leq k<14\\
                       \operatorname{bit}_k(m)=1}}D_{kj}.
\]
A graph-specific finite certificate assigns each mask either to
zero, to an adjacency-row prefix, or to a coordinate whose integer
numerator is neither 0 nor 3. Lean checks every certificate entry.
A structural proof of binary encoding ensures that every Boolean
prefix is represented, and reconstruction uniqueness identifies
the accepted vectors.

The final counterexample declarations are:
\begin{Verbatim}[fontsize=\small]
ACK.Witness550.isCounterexample
ACK.Witness563.isCounterexample
ACK.Witness570.isCounterexample
\end{Verbatim}
Each namespace also contains
\texttt{conjecture\_false}, proving the negation of
\texttt{ACK.Conjecture}.
The axiom reports for all six declarations contain exactly
\texttt{propext}, \texttt{Classical.choice}, and
\texttt{Quot.sound}. They contain neither
\texttt{sorryAx} nor \texttt{Lean.ofReduceBool}.
The formal verification requires no external SAT solver or
Python execution.

The current formalization proves the counterexample property and
the refutation. The rank argument, the additional graph invariants,
and the blowup corollary are not separately formalized.
\section*{An infinite family}
Replace each vertex $i$ by a nonempty independent set of size $t_i$ and each edge by the complete bipartite graph between its endpoint sets. Let $T$ be the resulting membership matrix, so $A'=TAT^T$. Since every set is nonempty, $T$ has a left inverse $L$.
Thus $AT^T=LA'$, while $A'=T(AT^T)$, so these two matrices
have the same row space. Consequently,
\[
 \operatorname{Row}(A')=\{xT^T:x\in\operatorname{Row}(A)\}.
\]
The vector $xT^T$ is Boolean if and only if $x$ is Boolean. Every nonzero such $x$ is a row of $A$, and its repeated-coordinate image is a row of $A'$. Hence every blowup is again an ACK counterexample, with rank 14. Taking one $t_i$ arbitrary and all others 1 yields a connected counterexample of every order at least 52. This is a corollary of the same refutation, not a separate conjecture settlement.
\section*{Reproduction}
The repository contains the explicit graphs, all integer
reconstruction and classification certificates as Lean source,
and the complete formal proofs.
The toolchain is Lean 4.35.0, with Mathlib pinned to revision
\texttt{56a005855ea17e0b0958f73ec801ae140b0a1679}.
The dependency manifest is committed.

After installing Lean through elan and obtaining repository access,
run the following commands from its root:
\begin{Verbatim}[fontsize=\small]
lake exe cache get
lake build
lake env lean Verification.lean
\end{Verbatim}
The final command displays the formal statement, the three
counterexample theorem types, and the axiom dependencies of all
three counterexample and refutation theorems. The GitHub workflow
also builds the project and runs this verification entry point.
\begin{thebibliography}{9}
\bibitem{ack} S. Akbari, P. J. Cameron, G. B. Khosrovshahi. \emph{Ranks and signatures of adjacency matrices}. Preprint (2004), Question 2. \url{https://webspace.maths.qmul.ac.uk/p.j.cameron/preprints/ranksign.pdf}
\bibitem{sciriha} I. Sciriha, B. Kaur, J. L. Borg, M. Debono. \emph{Potential counter-examples to a conjecture on the column space of the adjacency matrix}. Special Matrices 13 (2025), 20250033. \url{https://doi.org/10.1515/spma-2025-0033}.
\bibitem{akansha} S. Akansha, K. C. Sivakumar. \emph{Affirmative Results on a Conjecture on the Column Space of the Adjacency Matrix}. arXiv:2601.02746v1, 6 January 2026, Conjecture 1.1 and reference [6]. \url{https://arxiv.org/abs/2601.02746}
\end{thebibliography}

\end{document}
